% from ega4-10-fr.tex
\oldpage[IV-3]{102}
\begin{proposition}[10.4.2]
\label{IV.10.4.2-fr}
Soient $X$ un préschéma, $(U_\alpha)$ un recouvrement de $X$ formé d'ouverts
affines. Pour que $X$ soit un préschéma de Jacobson, il faut et il suffit que
les anneaux $\Gamma(U_\alpha,\mathcal O_X)$ des $U_\alpha$ soient des anneaux
de Jacobson.
\end{proposition}

Cela résulte de (10.3.3) et de la définition (10.4.1).

\begin{env}[10.4.3]
\label{IV.10.4.3-fr}
Un préschéma discret est un préschéma de Jacobson~; un anneau artinien est donc
un anneau de Jacobson. Tout anneau principal ayant une infinité d'idéaux
maximaux (par exemple $\mathbf Z$) est un anneau de Jacobson~; un anneau local
noethérien $A$ est un anneau de Jacobson si et seulement si son idéal maximal
est son seul idéal premier, c'est-à-dire si $A$ est artinien. Tout
sous-préschéma d'un préschéma de Jacobson est un préschéma de Jacobson en vertu
de (10.3.2).
\end{env}

\begin{proposition}[10.4.4]
\label{IV.10.4.4-fr}
Soit $B$ un anneau intègre. Les conditions suivantes sont équivalentes~:
\begin{enumerate}[label=\emph{\alph*)}]
  \item Il existe $f\neq0$ dans $B$ tel que $B_f$ soit un corps.
  \item Le corps des fractions de $B$ est une $B$-algèbre de type fini.
  \item Il existe un corps $K$ contenant $B$ et qui est une $B$-algèbre de
  type fini.
  \item Le point générique de $\operatorname{Spec}(B)$ est isolé dans
  $\operatorname{Spec}(B)$.
\end{enumerate}
\end{proposition}

Il est clair que \emph{d)} est équivalent à \emph{a)}, puisque \emph{d)}
signifie qu'il existe $f\neq0$ dans $B$ tel que $D(f)$ soit réduit au point
générique de $\operatorname{Spec}(B)$. Il est trivial que \emph{a)} entraîne
\emph{b)} et que \emph{b)} entraîne \emph{c)}. Enfin, \emph{c)} entraîne
\emph{a)}, en vertu de Bourbaki, \emph{Alg. comm.}, chap.~V, \S~3,
n\textsuperscript{o}~1, cor.~2 du th.~1.

\begin{proposition}[10.4.5]
\label{IV.10.4.5-fr}
Soit $A$ un anneau. Les conditions suivantes sont équivalentes~:
\begin{enumerate}[label=\emph{\alph*)}]
  \item $A$ est un anneau de Jacobson.
  \item Pour tout idéal premier non maximal $\mathfrak p$ de $A$ et tout
  $f\neq0$ dans $B=A/\mathfrak p$, $B_f$ n'est pas un corps.
  \item[\emph{b$'$)}] Toute $A$-algèbre de type fini $K$ qui est un corps est
  une $A$-algèbre finie.
\end{enumerate}
\end{proposition}

On sait que \emph{a)} entraîne \emph{b$'$)} (Bourbaki,
\emph{Alg. comm.}, chap.~V, \S~3, n\textsuperscript{o}~4, cor.~3 du th.~3).
En outre, le noyau de l'homomorphisme $A\to K$ est alors un idéal
\emph{maximal} $\mathfrak m$ de $A$, et $K$ est une extension \emph{finie} de
$A/\mathfrak m$ (\emph{loc. cit.}). Il est trivial que \emph{b$'$)} entraîne
\emph{b)}, puisque $A/\mathfrak p=B$ n'est pas un corps, toute $B$-algèbre de
type fini est une $A$-algèbre de type fini et $B_f$ est une $B$-algèbre de
type fini. Reste à voir que \emph{b)} implique \emph{a)}. Nous utiliserons le
lemme suivant~:

\begin{lemma}[10.4.5.1]
\label{IV.10.4.5.1-fr}
Soit $X$ un espace topologique ayant la propriété suivante~: pour toute partie
localement fermée $Z\neq\varnothing$ de $X$, il existe une partie $Z'$ de $Z$,
localement fermée dans $X$ (ou dans $Z$, ce qui revient au même), et un point
$x\in Z'$, fermé dans $Z'$. Alors, pour que $X$ soit un espace de Jacobson, il
faut et il suffit qu'aucun point non fermé $x$ de $X$ ne soit isolé dans
$\overline{\{x\}}$.
\end{lemma}

Si $X$ est un espace de Jacobson, un point non fermé $x$ de $X$ ne peut être
isolé dans $\overline{\{x\}}$, car cela signifierait qu'il existe un ouvert
$U$ de $X$ tel que $U\cap\overline{\{x\}}=\{x\}$~; or
$U\cap\overline{\{x\}}$ est localement fermé dans $X$ et ne contiendrait aucun
point fermé dans $X$, ce qui est contraire à l'hypothèse que $X$ est un espace
de Jacobson. Inversement, supposons vérifiée la condition de l'énoncé~; alors
l'ensemble des points fermés de $X$ est identique à l'ensemble $X_0$ des
$x\in X$ qui sont isolés dans $\overline{\{x\}}$~; mais il résulte de
(5.1.10.1) que

